\section{Sujet n\textsuperscript{o}\,22}
\noindent\begin{minipage}{0.5\linewidth}\footnotesize Office du Baccalauréat du Cameroun\end{minipage}%
\hfill\begin{minipage}{0.45\linewidth}\footnotesize\raggedleft Examen : \textbf{Baccalauréat} · Série : \textbf{C/E}\\
Épreuve : \textbf{Mathématiques} · Durée : \textbf{4\,h} · Coef. : \textbf{7}\end{minipage}
\par\smallskip{\color{onbo}\rule{\linewidth}{1pt}}

\subsection*{Partie A — Évaluation des ressources \hfill (15 points)}

\noindent\textbf{Exercice 1.} \hfill\textit{(3 points)}\\
Un cybercafé de Buea est équipé de machines provenant de trois lots. Le lot $A$ fournit
$50\,\%$ des machines, le lot $B$ en fournit $30\,\%$ et le lot $C$ les $20\,\%$ restants.
La probabilité qu'une machine tombe en panne au cours d'une journée est de $0{,}02$ pour
le lot $A$, $0{,}05$ pour le lot $B$ et $0{,}10$ pour le lot $C$.
\begin{enumerate}[label=\arabic*)]
\item On choisit une machine au hasard. Démontrer que la probabilité qu'elle tombe en
panne au cours d'une journée est $0{,}045$. \hfill\textit{(1,5 pt)}
\item Une machine est tombée en panne. Calculer la probabilité qu'elle provienne du
lot $C$. \hfill\textit{(1 pt)}
\item Le gérant choisit $3$ machines au hasard et de façon indépendante. Calculer la
probabilité qu'au moins une tombe en panne au cours de la journée. \hfill\textit{(0,5 pt)}
\end{enumerate}

\smallskip
\noindent\textbf{Exercice 2.} \hfill\textit{(3 points)}\\
Pour relier ses postes, le gérant achète des câbles courts coûtant \SI{800}{F} l'unité et
des câbles longs coûtant \SI{1300}{F} l'unité.
\begin{enumerate}[label=\arabic*)]
\item Justifier, à l'aide de l'algorithme d'Euclide, que $800$ et $1300$ sont premiers
entre eux après simplification, puis déterminer un couple de Bézout pour
$8u+13v=1$. \hfill\textit{(1,5 pt)}
\item En déduire toutes les solutions $(x\,;\,y)\in\Z^2$ de l'équation
$8x+13y=1$, puis résoudre dans $\N^2$ l'équation $8x+13y=200$. \hfill\textit{(1,5 pt)}
\end{enumerate}

\smallskip
\noindent\textbf{Exercice 3.} \hfill\textit{(4 points)}\\
On modélise le rendement (en milliers de francs) d'un poste de travail en fonction du
nombre $x$ d'heures d'utilisation par $f(x)=x-x\ln x$ pour $x>0$, de courbe $(\mathcal{C})$.
\begin{enumerate}[label=\arabic*)]
\item Déterminer $\displaystyle\lim_{x\to0^+}f(x)$ et $\displaystyle\lim_{x\to+\infty}f(x)$. \hfill\textit{(1 pt)}
\item Montrer que $f'(x)=-\ln x$, puis dresser le tableau de variations de $f$. \hfill\textit{(1,5 pt)}
\item À l'aide d'une intégration par parties, calculer $\displaystyle\int_1^{\mathrm{e}}x\ln x\,\dd x$,
puis en déduire $\displaystyle\int_1^{\mathrm{e}}f(x)\,\dd x$. \hfill\textit{(1,5 pt)}
\end{enumerate}

\smallskip
\noindent\textbf{Exercice 4.} \hfill\textit{(5 points)}\\
Le plan complexe est muni d'un repère orthonormé direct $(O\,;\,\vec u,\vec v)$. On considère
les points $A$, $B$, $C$ d'affixes respectives $z_A=2$, $z_B=3+i\sqrt3$ et $z_C=1+i\sqrt3$.
\begin{enumerate}[label=\arabic*)]
\item Soit $r$ la rotation de centre $O$ et d'angle $\frac\pi3$. Donner son écriture
complexe et démontrer que $r(A)=C$. \hfill\textit{(1,5 pt)}
\item Calculer $\dfrac{z_B-z_A}{z_C-z_A}$. En déduire la nature du triangle $ABC$. \hfill\textit{(1,5 pt)}
\item Soit $D$ l'image de $C$ par la translation de vecteur $\vec{AB}$. Déterminer
l'affixe $z_D$ et démontrer que $ABDC$ est un parallélogramme. \hfill\textit{(1 pt)}
\item Démontrer que $ABDC$ est en réalité un losange. \hfill\textit{(1 pt)}
\end{enumerate}

\subsection*{Partie B — Évaluation des compétences \hfill (5 points)}

\noindent\textit{Madame Eposi gère un centre informatique à Buea. Elle te sollicite, en tant
qu'élève de Terminale~C, pour trois décisions. Elle veut d'abord aménager une salle
rectangulaire : elle dispose de \SI{60}{\metre} de cloison mobile pour séparer une salle de
formation rectangulaire \emph{adossée à un mur} (le côté longeant le mur ne se cloisonne pas),
et cherche l'aire maximale. Ses abonnements rapportent \SI{150000}{F} le premier mois, puis
augmentent de \SI{6}{\percent} chaque mois. Enfin, sur l'ensemble de ses machines, $4{,}5\,\%$
tombent en panne un jour donné ; un technicien en contrôle $5$ au hasard.}

\begin{enumerate}[label=\textbf{Tâche \arabic*.},leftmargin=2.4em]
\item Déterminer les dimensions de la salle d'\emph{aire maximale} qu'elle peut aménager,
ainsi que cette aire. \hfill\textit{(1,5 pt)}
\item Déterminer le mois à partir duquel la recette mensuelle dépassera \SI{250000}{F}. \hfill\textit{(1,5 pt)}
\item Déterminer la probabilité qu'\emph{au moins une} des $5$ machines contrôlées soit en
panne. \hfill\textit{(1,5 pt)}
\end{enumerate}
\par\smallskip{\footnotesize\itshape Présentation générale : 0,5 point.}

% ════════════════════════════════════════════════════
\corrige{du sujet 22}

\noindent\textbf{Partie A — Exercice 1.}
On note $A,B,C$ les événements « la machine provient du lot $A,B,C$ » et $P$ « elle tombe en panne ».
1) Formule des probabilités totales :
$P(P)=0{,}5\times0{,}02+0{,}3\times0{,}05+0{,}2\times0{,}10=0{,}01+0{,}015+0{,}02=\mathbf{0{,}045}$.
2) Bayes : $P_P(C)=\dfrac{P(C)P_C(P)}{P(P)}=\dfrac{0{,}2\times0{,}10}{0{,}045}=\dfrac{0{,}02}{0{,}045}=\dfrac{4}{9}\approx0{,}444$.
3) Pour $3$ machines indépendantes : $P(\text{au moins une})=1-(1-0{,}045)^3=1-0{,}955^3\approx1-0{,}871=\mathbf{0{,}129}$.

\noindent\textbf{Exercice 2.}
1) Euclide : $1300=800\times1+500$, $800=500\times1+300$, $500=300\times1+200$,
$300=200\times1+100$, $200=100\times2+0$ : $\mathrm{pgcd}(800,1300)=100$, donc après division par
$100$, $8\wedge13=1$. Remontée : $1=100/100$ avec $1=3\times13-5\times8$ (car $3\times13=39$,
$5\times8=40$), soit $8\times(-5)+13\times3=1$ : $\mathbf{(u,v)=(-5\,;\,3)}$.
2) Solution particulière $(x_0,y_0)=(-5\,;\,3)$. Par différence, $8(x+5)=-13(y-3)$ ;
$13\mid x+5$ (Gauss, $8\wedge13=1$), d'où $x=-5+13k,\ y=3-8k$, $k\in\Z$.
Pour $8x+13y=200=200\times1$ : $x=-1000+13k,\ y=600-8k$. On veut $x\ge0$ et $y\ge0$ :
$x\ge0\Rightarrow k\ge76{,}9\Rightarrow k\ge77$ ; $y\ge0\Rightarrow k\le75$. Aucun entier
ne convient : il n'y a \textbf{aucune solution dans $\N^2$} pour $200$.

\noindent\textbf{Exercice 3.}
1) $\lim_{x\to0^+}x\ln x=0$ donc $\lim_{x\to0^+}f(x)=0$. En $+\infty$,
$f(x)=x(1-\ln x)\to-\infty$ car $1-\ln x\to-\infty$.
2) $f'(x)=1-(\ln x+x\cdot\frac1x)=1-\ln x-1=-\ln x$. Donc $f'(x)>0$ sur $]0,1[$ et
$f'(x)<0$ sur $]1,+\infty[$ : $f$ croît sur $]0,1]$, décroît sur $[1,+\infty[$, maximum $f(1)=1$.
3) IPP ($u=\ln x$, $v'=x$) : $\int_1^{\mathrm e}x\ln x\dd x=\big[\frac{x^2}{2}\ln x\big]_1^{\mathrm e}-\int_1^{\mathrm e}\frac{x}{2}\dd x
=\frac{\mathrm e^2}{2}-\big[\frac{x^2}{4}\big]_1^{\mathrm e}=\frac{\mathrm e^2}{2}-\frac{\mathrm e^2-1}{4}=\frac{\mathrm e^2+1}{4}$.
Puis $\int_1^{\mathrm e}f=\int_1^{\mathrm e}x\dd x-\int_1^{\mathrm e}x\ln x\dd x=\frac{\mathrm e^2-1}{2}-\frac{\mathrm e^2+1}{4}=\frac{\mathbf{\mathrm e^2-3}}{\mathbf 4}\approx1{,}10$.

\noindent\textbf{Exercice 4.}
1) $r(z)=\mathrm e^{i\pi/3}z=(\frac12+i\frac{\sqrt3}{2})z$. $r(A)$ a pour affixe
$(\frac12+i\frac{\sqrt3}{2})\times2=1+i\sqrt3=z_C$, donc $r(A)=C$.
2) $\frac{z_B-z_A}{z_C-z_A}=\frac{(3+i\sqrt3)-2}{(1+i\sqrt3)-2}=\frac{1+i\sqrt3}{-1+i\sqrt3}
=\frac{(1+i\sqrt3)(-1-i\sqrt3)}{(-1)^2+3}=\frac{-1-i\sqrt3-i\sqrt3+3}{4}=\frac{2-2i\sqrt3}{4}=\frac{1-i\sqrt3}{2}=\mathrm e^{-i\pi/3}$.
Le module vaut $1$ et l'argument $-\frac\pi3$ : $AB=AC$ et $(\vec{AC},\vec{AB})=-\frac\pi3$,
donc $ABC$ est \textbf{équilatéral}.
3) $z_D=z_C+(z_B-z_A)=(1+i\sqrt3)+(1+i\sqrt3)=2+2i\sqrt3$. Comme $\vec{CD}=\vec{AB}$
par construction, $ABDC$ est un parallélogramme.
4) $AB=|z_B-z_A|=|1+i\sqrt3|=2$ et $AC=|z_C-z_A|=|-1+i\sqrt3|=2$ : deux côtés consécutifs
égaux dans un parallélogramme $\Rightarrow$ \textbf{losange}.

\smallskip
\noindent\textbf{Partie B — Situation.}
\emph{Tâche 1.} Largeur $x$ (deux côtés), longueur $y=60-2x$. Aire $A(x)=x(60-2x)$ ;
$A'(x)=60-4x=0$ en $x=15$. Dimensions : $x=\SI{15}{\metre}$, $y=\SI{30}{\metre}$ ; aire
maximale $\SI{450}{\metre\squared}$.

\emph{Tâche 2.} Recette du mois $n$ : $u_n=150000\times1{,}06^{\,n-1}$. On résout $u_n>250000$,
soit $1{,}06^{\,n-1}>\frac53$, d'où $n-1>\frac{\ln(5/3)}{\ln1{,}06}\approx8{,}77$ : à partir du
\textbf{10\textsuperscript{e} mois} ($u_{10}\approx\SI{253400}{F}$).

\emph{Tâche 3.} $P(\text{au moins une})=1-(0{,}955)^5\approx1-0{,}7937=\mathbf{0{,}206}$, soit
environ \SI{20,6}{\percent}.
